
\begin{obeassessment}
\textbf{Kuis Bab 12}
1. Jelaskan perbedaan mendasar antara masalah logaritma diskrit dan masalah pemfaktoran bilangan bulat.
2. Uraikan langkah-langkah pembangkitan kunci pada sistem ElGamal.
3. Mengapa ElGamal disebut sebagai sistem enkripsi probabilistik?
\end{obeassessment}
\begin{competencychecklist}
\begin{tabularx}{\linewidth}{|X|c|}
\hline
Indikator Kompetensi & Tercapai \\
\hline
Saya dapat menjelaskan konsep logaritma diskrit dan akar primitif & [ ] \\
\hline
Saya dapat menghitung kunci publik dan privat ElGamal & [ ] \\
\hline
Saya dapat melakukan proses enkripsi dan dekripsi ElGamal & [ ] \\
\hline
Saya dapat menganalisis efek penggunaan kunci efemeral $k$ & [ ] \\
\hline
Saya dapat membandingkan ElGamal dengan RSA & [ ] \\
\hline
\end{tabularx}
\end{competencychecklist}
